\section{A stationary flow for noisy means}\label{sec:centering}

The probability transport turns conditional means into samples.  We now
construct the complementary operation: a path integral whose law is
$N(u,s^2I_d)$ for a prescribed $s>0$.  The key is to represent the centered
observable $M_*-u$ as a matrix adjoint divergence.  That matrix couples a
position with law $\xi_*$ to a Gaussian momentum while preserving their
product law; integrating the momentum equation then produces the noisy
mean.  This is still an exact mathematical construction: both its velocity
and its ideal initial law will need to be implemented later.

Retain the primitive and interpolation notation of
Sections~\ref{sec:preliminaries}--\ref{sec:divergence}.  Thus
$L=\lambda r$, $l=\lambda r^2\le1/2$, $M_0=u$, and the symmetric first
jet satisfies $DM_\rho=\rho U_1$.  Fix
\[
 0<R<1/2,\qquad b_*=(1-R^2)^{1/2},
\]
and use a star for evaluation at $\rho=b_*$.  The positive smoothing
radius $R$ keeps all coefficients smooth, although $F$ itself is only
$C^2$.  In this section $\rho$ is the Gaussian-interpolation parameter;
the new stationary flow will use a separate time variable $t$.

\subsection{Transport Jacobians and the zero endpoint}

Let $T_\rho=T_{\rho\to b_*}$ denote the probability transport
\eqref{eq:transport} from time $\rho$ to time $b_*$.  Given a terminal
position $y$, define
\[
 Y_\rho(y)=T_\rho^{-1}(y),\qquad
 B_\rho(y)=DT_\rho(Y_\rho(y))=(DY_\rho(y))^{-1}.
\]
In particular, $Y_{b_*}(y)=y$.  The backward path satisfies
$Y_\rho'=-rM_\rho(Y_\rho)$ with this terminal condition.  Its spatial
Jacobian $J_\rho=DY_\rho$ solves
\[
 J_\rho'=-rDM_\rho(Y_\rho)J_\rho,\qquad J_{b_*}=I_d.
\]
Differentiating its inverse gives the order of multiplication in
\begin{equation}\label{eq:backward-jacobian}
 B_\rho'=rB_\rho DM_\rho(Y_\rho)
         =r\rho B_\rho U_1(\rho,Y_\rho),\qquad B_0'=0.
\end{equation}
The last equality is an exact endpoint identity, not an asymptotic
estimate.  Also $\|B_\rho\|_{\op}\le C$, because
$r\|DM_\rho\|_{\op}\le Cl$ and the time interval has length at most
one.  The following estimates control both spatial differentiation of
the kernel and time interpolation of $B_\rho$.

\begin{lemma}[Transport derivatives]\label{lem:transport-derivatives}
For each fixed $k\ge1$, the order-$k$ spatial derivatives of any transport
between two times in $[0,b_*]$, after subtracting the identity map, have
all nontrivial split norms at most
\[
 C_k lR^{-(k-1)}.
\]
For each fixed $m\ge1$ and fixed finite $p\ge1$, uniformly in
$0\le\rho\le b_*$,
\[
 \big\|\|B_\rho^{(m)}(y)\|_{\op}\big\|_{L^p(y\sim\xi_*)}
 \le \polylog(d)\,l(1-\rho^2)^{-m}.
\]
The constants and the polylogarithmic factor may depend on the fixed
orders and on $p$, but not on the primitive center or scale.
\end{lemma}

\begin{proof}
Differentiate the transport equation with respect to its starting
position.  Gronwall's inequality gives a bounded first derivative and an
$O(l)$ bound for its difference from $I_d$.  For order $k>1$, the chain
rule has terms consisting of an order-$j$ derivative of the velocity
and derivatives of the state of orders $i_1,\ldots,i_j\ge1$, where
$i_1+\cdots+i_j=k$.  By \eqref{eq:jets}, the outer tensor has every
split norm at most $C_j lR^{-(j-1)}$.  Attach the inner tensors one at
a time.  Each attachment is a single contraction retaining the outer
output and the inner derivative slots, so
Lemma~\ref{lem:split-contraction} applies without a dimension factor.
Bound a first state derivative by a constant and use the induction
hypothesis for all higher state derivatives.  Since $l\le1$, the total
power of $R^{-1}$ is at most
\[
 (j-1)+\sum_{\nu=1}^j(i_\nu-1)=k-1,
\]
and each forcing term retains the outer factor $l$.  The term involving
the unknown order-$k$ state derivative is linear, with coefficient
bounded by $Cl$, and this derivative starts at zero.  Gronwall proves
the assertion.  The same argument with absolute interval lengths
applies to transport in the reverse direction.

For the second assertion, write
$C_\rho=rDM_\rho(Y_\rho)=r\rho U_1(\rho,Y_\rho)$, so that
$B_\rho'=B_\rho C_\rho$.  The product rule expresses $B_\rho^{(m)}$
as a finite sum of ordered products of $B_\rho$ and $k\ge1$ factors
$C_\rho^{(j_1)},\ldots,C_\rho^{(j_k)}$, with
$j_1+\cdots+j_k=m-k$.  The zero-order jet bound and
\eqref{eq:path-jet} imply, for every fixed moment order,
\[
 \big\|\|C_\rho^{(j)}\|_{\op}\big\|_{L^p}
 \le \polylog(d)\,l(1-\rho^2)^{-j}.
\]
Here the path estimates apply because inverse transport sends
$y\sim\xi_*$ to $Y_\rho(y)\sim\xi_\rho$.  Apply H\"older with
sufficiently high fixed moments to each product.  Using $l^k\le l$
and $(1-\rho^2)^{-(m-k)}\le(1-\rho^2)^{-m}$ gives the claimed
conservative exponent.  Moments below one, if desired, follow from
the $L^1$ bound.
\end{proof}

\subsection{A kernel for the centered observable}

For a matrix field $K$, define its row-wise adjoint divergence by
\[
 (\nabla_{\xi_*}^*K)_i
 =\sum_j\{-\partial_jK_{ij}+(\partial_jH_*)K_{ij}\}.
\]
Thus $\nabla_{\xi_*}^*K=h$ means
$\E_{\xi_*}[h_i\phi]=\E_{\xi_*}\sum_jK_{ij}\partial_j\phi$
for smooth compactly supported scalar tests $\phi$.  We seek this
identity with $h=M_*-u$.

This is a Stein-type integration-by-parts identity.  A classical
coordinate Stein kernel for a law $\xi$ represents the observable
$y-\E_\xi y$ instead.  Courtade, Fathi and Pananjady construct such
kernels from spectral-gap assumptions
\citep[arXiv version~2, Definition~2.1 and Theorem~2.4]{courtade2019stein}.
Our observable is $M_*-u$; its mean is zero by conditional expectation.
The transport formula below supplies its particular adjoint-divergence
representation together with pointwise derivative bounds.  Those bounds
are needed for quadrature and do not follow from the cited existence
statement.  We do not require $K$ to be symmetric or positive.

Define
\begin{equation}\label{eq:stein-kernel}
 K(y)^T=\int_0^{b_*}B_\rho(y)U_1(\rho,Y_\rho(y))\,d\rho
       =\frac1r\int_0^{b_*}B_\rho'(y)\,\frac{d\rho}{\rho}.
\end{equation}
The first integral defines the integrand at zero in the second formula:
$B_\rho'/(r\rho)=B_\rho U_1(\rho,Y_\rho)$ extends smoothly there.
The transpose is essential: the adjoint divergence will act on the
rows of $K$, while the position velocity will use $K^T$.

\begin{lemma}[The centering kernel]\label{lem:centering-kernel}
The field $K$ is smooth and satisfies
\[
 \nabla_{\xi_*}^*K=M_*-u,\qquad \|K(y)\|_{\op}\le CL.
\]
For each fixed $j\ge1$, every nontrivial split norm of $\nabla^jK$ is
bounded pointwise by $C_jLR^{-j}$.
\end{lemma}

\begin{proof}
Fix an output coordinate $i$ and a smooth compactly supported scalar
function $\phi$.  The transported test $\phi\circ T_\rho$ has zero
material derivative: it is constant along each probability-transport
path.  Differentiate its product with $M_{\rho,i}$ under the
transported laws.  Lemma~\ref{lem:material-identities} gives
$\mathcal DM_\rho=\nabla_{\xi_\rho}^*U_1$, so integration by parts
and the chain rule yield
\[
 \frac{d}{d\rho}\E_{\xi_\rho}[M_{\rho,i}\,\phi\circ T_\rho]
 =\E_{\xi_\rho}\sum_{j,k}(U_1)_{ik}(DT_\rho)_{jk}
                                  (\partial_j\phi)\circ T_\rho.
\]
Smooth transport on the compact time interval and Gaussian tails
justify differentiation; alternatively the test can first be cut off
and the resulting identities passed to the limit.  Push the expectation
forward to $y\sim\xi_*$.  Since $U_1$ is symmetric, its coefficient of
$\partial_j\phi(y)$ is
\[
 \sum_k (U_1)_{ik}(\rho,Y_\rho(y))(B_\rho(y))_{jk}
       =(B_\rho(y)U_1(\rho,Y_\rho(y)))_{ji}.
\]
The terminal expectation is $\E_{\xi_*}[M_{*,i}\phi]$.
At zero, $M_0=u$ and $(T_0)_\#\xi_0=\xi_*$, so the initial
expectation is $u_i\E_{\xi_*}\phi$.  Integrating the displayed
derivative in $\rho$ proves
\[
 \E_{\xi_*}[(M_{*,i}-u_i)\phi]
   =\E_{\xi_*}\sum_jK_{ij}\partial_j\phi.
\]
This proves the adjoint-divergence identity weakly, and smoothness
makes it a pointwise identity.

For the bounds, $\|B_\rho\|_{\op}\le C$ and
$\|U_1\|_{\op}\le CL$ immediately give $\|K\|_{\op}\le CL$.
For a positive-order spatial derivative, use
Lemma~\ref{lem:transport-derivatives} on both $DT_\rho$ and
$Y_\rho$.  An order-$j$ derivative of $DT_\rho$ has all split norms
at most $C_jlR^{-j}$, and an order-$j$ derivative of $U_1$ has all
split norms at most $C_jLR^{-j}$ by \eqref{eq:jets}.  When these
fields are composed with $Y_\rho$, attach each chain-rule tensor by
a single edge as in the preceding lemma.  First derivatives of
$Y_\rho$ are bounded, and a derivative of order $i>1$ costs
$C_i lR^{-(i-1)}$.  Every order-$j$ term in the product has total
cost at most $C_jLR^{-j}$, since $l\le1$.  Each factor retains free
slots, so the single-contraction bound preserves every nontrivial
split.  Integrating over an interval of length at most one proves
the assertion.  All orders are smooth by Gaussian conditional
integration and the smooth transport on $\rho\le b_*<1$; this uses
no higher derivatives of the original unsmoothed gradient.
\end{proof}

\subsection{Stationarity and the path integral}

The following dynamics specialize the skew-matrix construction of
Ma, Chen and Fox
\citep[Section~2.1, Equation~(3), Theorem~1]{ma2015recipe}.
The transport-derived $K$ makes the momentum displacement encode
$M_*-u$.  We prove completeness and invariance directly, since a
formal stationary-density calculation alone would not ensure a
solution exists for the full integration interval.

\begin{proposition}[A noisy mean from a path integral]
\label{prop:centering-law}
For $s>0$, let $(Y_0,G_0)$ have law $\xi_*\otimes N(0,I_d)$ and solve
\begin{equation}\label{eq:centering-flow}
 \dot Y_t=-K(Y_t)^TG_t/s,\qquad
 \dot G_t=(M_*(Y_t)-u)/s.
\end{equation}
The solution exists for every finite time, and its law remains
$\xi_*\otimes N(0,I_d)$.  In particular,
\begin{equation}\label{eq:centering-output}
 sG_0+\int_0^1M_*(Y_t)\,dt\ \stackrel{\mathrm{law}}=
 u+sG,\qquad G\sim N(0,I_d).
\end{equation}
\end{proposition}

\begin{proof}
The vector field $v$ in \eqref{eq:centering-flow} is smooth and hence
locally Lipschitz.  The bound on $K$ and the global Lipschitz bound on
$M_*$ imply $|v(y,g)|\le C_{x,r,s}(1+|y|+|g|)$.  Gronwall's
inequality therefore prevents a solution from escaping to infinity
in finite positive or negative time.  Local uniqueness gives a
complete smooth flow $\Phi_t$, with inverse $\Phi_{-t}$.  We only
need local Lipschitz continuity here; the factor $G$ does not in
general give a globally Lipschitz joint vector field.

Write the product density as
\[
 p(y,g)=(2\pi)^{-d/2}\exp(-H_*(y)-|g|^2/2).
\]
Differentiating in the position and momentum variables separately,
with the row convention of the kernel lemma, gives
\[
 \operatorname{div}_y(pv_Y)
   =\frac{p}{s}\sum_i g_i(\nabla_{\xi_*}^*K)_i
   =\frac{p}{s}\,g\cdot(M_*-u),\qquad
 \operatorname{div}_g(pv_G)
   =-\frac{p}{s}\,g\cdot(M_*-u).
\]
Thus $\operatorname{div}(pv)=0$.  The Jacobian formula for the
complete flow now yields, pointwise,
\[
 \frac{d}{dt}\{p(\Phi_t(z))\det D\Phi_t(z)\}
   =\operatorname{div}(pv)(\Phi_t(z))\det D\Phi_t(z)=0.
\]
The determinant starts at one and stays positive.  Change of variables
therefore proves that $\Phi_t$ preserves the probability density $p$.

For later differentiation it is useful to record the same cancellation
as an operator identity.  Relative to the product law, the Lie
derivative is
\begin{equation}\label{eq:centering-skew-operator}
 \mathcal L=\nabla^*A\nabla,\qquad
 A=\begin{pmatrix}0&K^T/s\\-K/s&0\end{pmatrix}.
\end{equation}
Indeed the second-order part vanishes by skew symmetry.  The
first-order coefficients are $-K^Tg/s$ in position and
$\nabla_{\xi_*}^*K/s=(M_*-u)/s$ in momentum, exactly as in
\eqref{eq:centering-flow}.

Finally, integrating the momentum equation over $[0,1]$ gives the
pathwise identity
\[
 sG_0+\int_0^1M_*(Y_t)\,dt=u+sG_1.
\]
Stationarity makes $G_1$ a standard Gaussian, proving
\eqref{eq:centering-output}.
\end{proof}

Stationarity also makes $G_1$ independent of $Y_1$.  It does not assert
that $G_1$ is independent of the initial anchors $(Y_0,G_0)$.
The exact mean $u=M_0$ remains in the momentum velocity, and the
proposition assumes an ideal sample from $\xi_*$.  Thus the proposition
is a comparison construction, not yet an oracle routine.  Section~\ref{sec:numerics}
will express its approximation by finitely many exact mean evaluations;
Section~\ref{sec:compiler} will replace every such evaluation, including
$M_0$.  The displayed output itself needs no separate addition of $u$.

\subsection{Derivatives on a stationary trajectory}

To approximate the path integral by high-order quadrature, we need
bounds on time derivatives at each fixed time.  Stationarity allows
these bounds to be proved in the fixed product measure.  The following
lemma separates that argument from the particular kernel and avoids
requiring a bound on the supremum of a random path.

\begin{lemma}[Stationary time derivatives]\label{lem:stationary-derivatives}
Fix $0<R\le1$.  Let a smooth probability density on $\R^n$, with
$n\le Cd$, have Gaussian tails and negative logarithm $H$.  Assume
its Hessian is bounded and that every positive-order score derivative
$\nabla^j\nabla H$ has every one-slot split norm at most
$C_jR^{-(j-1)}$, for $j\ge1$.  Here a one-slot split isolates one
tensor slot as the matrix input and groups all remaining slots as
its output.  Suppose a stationary autonomous flow has Lie derivative
$\mathcal L=\nabla^*A\nabla$, where the adjoint is for this density,
$A$ is skew symmetric, and
\[
 |\nabla^jA|_{\text{one-slot split}}\le C_j\alpha R^{-j}
 \quad(j\ge0)
\]
for every one-slot split.  Let $\phi_0$ be smooth and scalar- or
vector-valued, with
\[
 |\nabla^j\phi_0|_F\le C_ja_0\sqrt d\,R^{-j}
 \quad(j\ge1).
\]
Assume the fields and all derivatives used have polynomial growth.
Then, for each fixed $k\ge1$, the order-$k$ time derivative of
$\phi_0$ on the stationary flow has RMS norm at most
\begin{equation}\label{eq:stationary-derivatives}
 C_k a_0\sqrt d\,(\alpha R^{-2})^k.
\end{equation}
\end{lemma}

\begin{proof}
Put $\phi_k=\mathcal L^k\phi_0$.  We prove, for $j\ge0$ and
$k\ge1$, the stronger estimate
\begin{equation}\label{eq:stationary-spatial-induction}
 \|\nabla^j\phi_k\|_{F,2}
 \le C_{j,k}a_0\sqrt d\,\alpha^kR^{-(j+2k)}.
\end{equation}
For $k=0$ the same bound is assumed only for $j\ge1$, which suffices:
the iteration differentiates $\phi_{k-1}$ before multiplying by $A$.
All $L^2$ norms below are for the invariant density.

Set $V=A\nabla\phi_{k-1}$.  Leibniz's rule and the induction
hypothesis give, for each $h\ge0$,
\[
 \|\nabla^h V\|_{F,2}
 \le C_{h,k}a_0\sqrt d\,\alpha^k
                         R^{-(h+1+2(k-1))}.
\]
To see why no dimension factor appears, in each coefficient tensor
$\nabla^vA$ isolate the slot contracted against the gradient of
$\phi_{k-1}$.  Its one-slot operator norm multiplies the Frobenius
norm of the other factor.  Summing the finitely many Leibniz terms
preserves the displayed estimate.

Now commute $j$ derivatives past the adjoint divergence in
$\phi_k=\nabla^*V$.  The main term is $\nabla^*(\nabla^jV)$.
The Frobenius estimate \eqref{eq:divergence-l2} bounds its $L^2$ norm
by a constant times
$\|\nabla^jV\|_{F,2}+\|\nabla^{j+1}V\|_{F,2}$.
These two terms have powers $R^{-(j+2k-1)}$ and
$R^{-(j+2k)}$, respectively.  Each remaining commutator contracts
an order-$v$ derivative of the score, $1\le v\le j$, with
$\nabla^{j-v}V$.  The coefficient's one-slot norm again suffices;
the total exponent of $R^{-1}$ is
\[
 (v-1)+(j-v+1+2(k-1))=j+2k-2.
\]
All these terms are bounded by
\eqref{eq:stationary-spatial-induction} because $0<R\le1$.
This closes the induction, including $j=0$ for $k\ge1$.
Gaussian tails and polynomial growth justify the integrations by
parts through cutoffs.

Along an autonomous trajectory, the $k$th time derivative of
$\phi_0$ is $\phi_k$ evaluated at the current state.  Stationarity
identifies its RMS norm at every time with the fixed-measure norm
in \eqref{eq:stationary-spatial-induction} for $j=0$.
\end{proof}

For \eqref{eq:centering-flow}, the ambient dimension is $2d$, and
\eqref{eq:centering-skew-operator} satisfies the lemma's coefficient
bounds with $\alpha=L/s$.  The product potential is
$H_*(y)+|g|^2/2$; its score-derivative bounds follow from
\eqref{eq:jets} and the Gaussian factor.  For the joint state take
$a_0=1$: its first derivative has Frobenius norm $\sqrt{2d}$ and
all higher derivatives vanish.  For $M_*$ take $a_0=L$.
Indeed, isolating any derivative slot in \eqref{eq:jets} gives a
matrix with input dimension $d$, so passage from its operator norm
to its Frobenius norm costs at most $\sqrt d$.
Consequently, at every stationary time,
\[
 \left\|\frac{d^k}{dt^k}(Y_t,G_t)\right\|_2
       \le C_k\sqrt d\,((L/s)R^{-2})^k,\qquad
 \left\|\frac{d^k}{dt^k}M_*(Y_t)\right\|_2
       \le C_kL\sqrt d\,((L/s)R^{-2})^k.
\]
The strict radius choice $0<R<1/2$ made for this construction
satisfies the abstract lemma's $0<R\le1$ hypothesis.

We also need an initial displacement estimate for the later iteration.
Conditional expectation and the primitive Poincar\'e inequality give
\[
 \E_{\xi_*}|M_*-u|^2
 \le\E_\nu|f-u|^2\le CL^2d.
\]
At every stationary time, $G_t$ is independent of $Y_t$, hence
\[
 \E|K(Y_t)^TG_t|^2
   =\E\|K(Y_t)\|_F^2\le CL^2d.
\]
Integrate the velocity and apply Minkowski's inequality to obtain
\begin{equation}\label{eq:centering-initial-displacement}
 \|(Y_t,G_t)-(Y_0,G_0)\|_2
 \le Ct(L/s)\sqrt d\le C(L/s)\sqrt d
 \qquad(0\le t\le1).
\end{equation}
These are fixed-time RMS estimates for the exact stationary flow.
They give the derivative and starting-error bounds used in the next
section to approximate its path integral by finite formulas.
